\chapter{Preliminaries}
\label{chap:preliminaries}
% Paper label sec:prelims kept for cross-reference fidelity; resolves to chapter number.
\label{sec:prelims}

\section{Polynomials and measurements}

For complex numbers $\alpha,\beta \in \C$, write $\alpha \approx_\eps \beta$ when $|\alpha-\beta|\le \eps$. In particular,
\begin{equation}\label{eq:triangle-inequality-for-numbers}
	\text{if $\alpha \approx_{\eps} \beta$ and $\beta \approx_{\delta} \gamma$, then $\alpha \approx_{\eps+\delta}\gamma$.}
\end{equation}

Fix a prime power $q=p^t$. Write $\omega=e^{2\pi i/p}$, and let $\tr:\F_q\to\F_p$ denote the finite-field trace.

\begin{definition}[Finite field trace]\label{def:ff-trace}
	\lean{MIPStarRE.LDT.Preliminaries.ffTrace}
	\leanok
	The finite-field trace is the map $\tr:\F_q\to\F_p$ defined by
	\[
		\tr[x] = \sum_{\ell=0}^{t-1} x^{p^\ell}.
	\]
\end{definition}

\begin{lemma}\label{prop:fourier-fact-scalar}
	\lean{MIPStarRE.LDT.Preliminaries.fourier_fact_scalar}
	\leanok
	\uses{def:ff-trace}
	Let $a \in \F_q$. Then
	\begin{equation*}
		\E_{\bx \sim \F_q} \omega^{\tr[\bx \cdot a]}
		= \left\{\begin{array}{rl}
			1 & \text{if $a = 0$}, \\
			0 & \text{otherwise}.
		\end{array}\right.
	\end{equation*}
\end{lemma}
\begin{proof}
	If $a=0$, then $\tr[\bx\cdot a]=0$ for every $\bx$, so the average is $1$. If $a\neq 0$, choose $y \in \F_q$ such that $\tr[a\cdot y]\neq 0$, and set
	\[
		C=\E_{\bx \sim \F_q} \omega^{\tr[\bx \cdot a]}.
	\]
	Translation invariance of the uniform distribution gives
	\[
		C=\E_{\bx \sim \F_q} \omega^{\tr[(\bx+y)\cdot a]}
		= C \cdot \omega^{\tr[y\cdot a]}.
	\]
	Since $\omega^{\tr[y\cdot a]}\neq 1$, it follows that $C=0$.
\end{proof}

\begin{lemma}\label{prop:fourier-fact-vector}
	\lean{MIPStarRE.LDT.Preliminaries.fourier_fact_vector}
	\leanok
	\uses{def:ff-trace, prop:fourier-fact-scalar}
	Let $v \in \F_q^m$. Then
	\begin{equation*}
		\E_{\bu \sim \F_q^m} \omega^{\tr[\bu \cdot v]}
		= \left\{\begin{array}{rl}
			1 & \text{if $v = 0$}, \\
			0 & \text{otherwise}.
		\end{array}\right.
	\end{equation*}
\end{lemma}
\begin{proof}
	By linearity of the trace and independence of the coordinates,
	\[
		\E_{\bu \sim \F_q^m} \omega^{\tr[\bu \cdot v]}
		= \prod_{i=1}^m \E_{\bu_i \sim \F_q} \omega^{\tr[\bu_i \cdot v_i]}.
	\]
	Proposition~\ref{prop:fourier-fact-scalar} shows that the $i$-th factor is $1$ when $v_i=0$ and $0$ otherwise, so the product is $1$ exactly when $v=0$.
\end{proof}

\begin{definition}[Low individual degree polynomials]\label{def:polyfunc}
	\lean{MIPStarRE.LDT.Preliminaries.polyFunc, MIPStarRE.LDT.Polynomial}
	Let $\polyfunc{m}{q}{d}$ be the set of polynomials $g \in \F_q[x_1,\dots,x_m]$ whose degree in each variable is at most $d$, viewed as functions $\F_q^m \to \F_q$ via evaluation. (Over finite fields, distinct low-degree polynomials can induce the same function; we identify elements of $\polyfunc{m}{q}{d}$ with their polynomial representatives, not their functional equivalence classes.)
\end{definition}

\begin{remark}\label{rem:individual-degree-convention}
	\lean{MIPStarRE.LDT.Preliminaries.polyFuncMonotone}
	Here, ``individual degree $d$'' means degree at most $d$ in each coordinate. In
	particular,
	\[
		\polyfunc{m}{q}{d} \subseteq \polyfunc{m}{q}{d+1}.
	\]
\end{remark}

\begin{lemma}[Schwartz--Zippel lemma~\cite{Sch80,Zip79}]\label{lem:schwartz-zippel-total-degree}
	\lean{MIPStarRE.LDT.Preliminaries.schwartzZippel_totalDegree}
	\leanok
	Let $g, h:\F_q^m \rightarrow \F_q$ be two distinct polynomials of total degree~$d$.
	Then
	\begin{equation*}
		\Prb_{\bx \sim \F_q^m}[g(\bx) = h(\bx)] \leq \frac{d}{q}.
	\end{equation*}
\end{lemma}
\begin{proof}
	Apply the Schwartz--Zippel lemma to the nonzero polynomial $g-h$, which has total degree at most $d$.
\end{proof}

\begin{lemma}[Schwartz--Zippel for individual degree]\label{lem:schwartz-zippel-individual}
	\lean{MIPStarRE.LDT.Preliminaries.schwartzZippel_individualDegree}
	\leanok
	\uses{def:polyfunc, lem:schwartz-zippel-total-degree}
	If $g,h \in \polyfunc{m}{q}{d}$ are distinct, then
	\[
		\Prb_{u \sim \F_q^m}[g(u)=h(u)] \le \frac{md}{q}.
	\]
\end{lemma}
\begin{proof}
	The difference $g-h$ is a nonzero polynomial of total degree at most $md$. Apply Lemma~\ref{lem:schwartz-zippel-total-degree}.
\end{proof}

\begin{definition}[Measurements and submeasurements]\label{def:submeasurement}
	\lean{MIPStarRE.LDT.SubMeas, MIPStarRE.LDT.Measurement, MIPStarRE.LDT.ProjSubMeas}
	\leanok
	Let $\mathcal H$ be a Hilbert space and let $\mathcal A$ be a set of outcomes. A submeasurement on $\mathcal A$ is a family $A=\{A_a\}_{a \in \mathcal A}$ of Hermitian positive semidefinite operators on $\mathcal H$ such that $\sum_a A_a \le I$. It is a measurement when $\sum_a A_a = I$, and it is projective when each $A_a$ is an idempotent projection. (In the Lean formalization, outcome sets are assumed finite throughout.)
\end{definition}

\begin{definition}[Low-degree polynomial measurements]\label{def:polymeasurements}
	\lean{MIPStarRE.LDT.SubMeas, MIPStarRE.LDT.Measurement, MIPStarRE.LDT.Polynomial}
	\leanok
	\uses{def:polyfunc, def:submeasurement}
	We write $\polysub{m}{q}{d}$ for the submeasurements indexed by $\polyfunc{m}{q}{d}$, and $\polymeas{m}{q}{d}$ for the corresponding measurements.
\end{definition}

\begin{definition}[Post-processing]\label{def:post-processing}
	\lean{MIPStarRE.LDT.postprocess}
	\leanok
	\uses{def:submeasurement}
	Let $A=\{A_a\}_{a \in \mathcal A}$ be a family of operators and let $f\colon \mathcal A \to \mathcal B$. The post-processed family $A_{[f(a)=b]}$ is defined by
	\[
		A_{[f(a)=b]} = \sum_{a\, :\, f(a)=b} A_a.
	\]
\end{definition}

\begin{proposition}[Post-processing preserves measurements]\label{prop:post-processing-preserves}
	\lean{MIPStarRE.LDT.Preliminaries.postprocessPreservesMeasurements}
	\leanok
	\uses{def:post-processing, def:submeasurement}
	Let $A=\{A_a\}_{a \in \mathcal A}$ be a family of operators and let
	$f\colon \mathcal A \to \mathcal B$. Then
	\[
		\sum_a A_a = \sum_b A_{[f(a)=b]}.
	\]
	Consequently, if $\{A_a\}$ is a submeasurement, respectively a measurement, then
	$\{A_{[f(a)=b]}\}$ is again a submeasurement, respectively a measurement.
\end{proposition}
\begin{proof}
	Each $a \in \mathcal A$ contributes exactly once to the right-hand side, namely in
	the summand indexed by $b=f(a)$.
\end{proof}

\begin{remark}[Post-processing notation]\label{rem:post-processing-notation}
	In the notation $A_{[f(a)=b]}$, the measurement outcome is the variable to which
	the function is applied. For example, if $G=\{G_g\} \in \polysub{m}{q}{d}$ and
	$u \in \F_q^m$, then
	\[
		G_{[g(u)=a]} = \sum_{g\,:\,g(u)=a} G_g.
	\]
	Here $g$ is the measurement outcome, and $u$ is the evaluation point.
\end{remark}

\begin{definition}[Completion of a submeasurement]\label{def:measurement-completion}
	\lean{MIPStarRE.LDT.Preliminaries.completeAtOutcome}
	\leanok
	\uses{def:submeasurement}
	Let $A=\{A_a^x\}_{a \in \mathcal A}$ be a submeasurement. Its completion is the measurement $\widehat A$ with outcome set $\widehat{\mathcal A}=\mathcal A \cup \{\bot\}$ given by $\widehat A_a^x=A_a^x$ for $a \in \mathcal A$ and $\widehat A_\bot^x = I-\sum_a A_a^x$.
\end{definition}

\section{Consistency and state-dependent distance}
\label{sec:comparing-measurements}


\begin{definition}[Consistency]\label{def:simeq}
	\lean{MIPStarRE.LDT.ConsRel}
	\leanok
	Let $\ket{\psi} \in \mathcal H_{\mathrm A} \ot \mathcal H_{\mathrm B}$. Let $A=\{A_a^x\}$ and $B=\{B_a^x\}$ be submeasurements with the same answer set, and let $\mathcal D$ be a distribution on the question set. We write
	\[
		A_a^x \ot I \simeq_\delta I \ot B_a^x
	\]
	when
	\begin{equation}\label{eq:no-big-Oh}
		\E_{x \sim \mathcal D} \sum_{a \ne b} \bra{\psi} A_a^x \ot B_b^x \ket{\psi} \le \delta.
	\end{equation}
\end{definition}

\begin{lemma}[Characterization of a good symmetric strategy]\label{lem:good-strategy-characterization}
	\lean{MIPStarRE.LDT.Preliminaries.goodStrategyCharacterization}
	\leanok
	\uses{def:good-strategy, def:simeq, prop:simeq-for-measurements}
	A symmetric projective strategy $(\psi,A,B,L)$ is $(\eps,\delta,\gamma)$-good if and only if the following hold on the relevant test distributions:
	\[
		A_a^u \ot I \simeq_\eps I \ot B^{\ell}_{[f(u)=a]},
		\qquad
		A_a^u \ot I \simeq_\delta I \ot A_a^u,
		\qquad
		A_a^u \ot I \simeq_\gamma I \ot L^{\ell}_{[f(u)=a]}.
	\]
\end{lemma}
\begin{proof}
	Each subtest accepts exactly when the two measurement outcomes agree. Because the relevant families are measurements, Lemma~\ref{prop:simeq-for-measurements} rewrites those acceptance probabilities as consistency bounds.
\end{proof}

\begin{lemma}[Consistency for measurements]\label{prop:simeq-for-measurements}
	\lean{MIPStarRE.LDT.Preliminaries.simeqForMeasurements}
	\leanok
	\uses{def:simeq}
	If $A=\{A_a^x\}$ and $B=\{B_a^x\}$ are measurements, then
	\[
		A_a^x \ot I \simeq_\delta I \ot B_a^x
		\qquad\text{if and only if}\qquad
		\E_x \sum_a \bra{\psi} A_a^x \ot B_a^x \ket{\psi} \ge 1-\delta.
	\]
\end{lemma}
\begin{proof}
	Expand the off-diagonal sum by writing $\sum_{a \ne b} A_a^x \ot B_b^x = \sum_b (I-A_b^x) \ot B_b^x$.
\end{proof}

\begin{definition}[State-dependent distance]\label{def:approx_delta}
	\lean{MIPStarRE.LDT.SDDRel}
	\leanok
	Let $\ket{\psi} \in \mathcal H$, let $A=\{A_a^x\}$ and $B=\{B_a^x\}$ be families of operators on $\mathcal H$, and let $\mathcal D$ be a distribution on the question set. We write
	\[
		A_a^x \approx_\delta B_a^x
	\]
	when
	\[
		\E_{x \sim \mathcal D} \sum_a \lVert (A_a^x-B_a^x)\ket{\psi} \rVert^2 \le \delta.
	\]
\end{definition}

In particular, if $A^x_a \ot I \simeq_{\eps} I \ot C^x_a$, the question is what hypothesis on $A$ and $B$ allows one to conclude that
\begin{equation}\label{eq:can-we-use-approx-delta-to-derive-this}
	B^x_a \ot I \simeq_{\eps} I \ot C^x_a.
\end{equation}

\begin{lemma}[Transferring consistency through state-dependent distance]\label{prop:triangle-sub}
	\lean{MIPStarRE.LDT.Preliminaries.triangleSub}
	\leanok
	\uses{def:simeq, def:approx_delta, prop:simeq-for-measurements}
	Let $A=\{A_a^x\}$ and $B=\{B_a^x\}$ be measurements and let $C=\{C_a^x\}$ be a submeasurement. If $A_a^x \ot I \simeq_\delta I \ot C_a^x$ and $A_a^x \ot I \approx_\eps B_a^x \ot I$, then
	\[
		B_a^x \ot I \simeq_{\delta+\sqrt\eps} I \ot C_a^x.
	\]
\end{lemma}
\begin{proof}
	\leanok
	Rewrite the inconsistency with $C$ as total mass minus diagonal overlap. The total mass is unchanged because $A$ and $B$ are measurements, and the diagonal overlap changes by at most $\sqrt\eps$ by Cauchy--Schwarz.
\end{proof}

\begin{lemma}[Consistency implies state-dependent distance for measurements]\label{prop:simeq-to-approx}
	\lean{MIPStarRE.LDT.Preliminaries.simeqToApprox}
	\leanok
	\uses{def:simeq, def:approx_delta, prop:simeq-for-measurements}
	If $A$ and $B$ are measurements and $A_a^x \ot I \simeq_\delta I \ot B_a^x$, then
	\[
		A_a^x \ot I \approx_{2\delta} I \ot B_a^x.
	\]
	If both measurements are projective, the converse also holds.
\end{lemma}
\begin{proof}
	Expand the squared norm of $(A_a^x \ot I - I \ot B_a^x)\ket{\psi}$ and use the diagonal-overlap formula from Lemma~\ref{prop:simeq-for-measurements}.
\end{proof}

\begin{remark}\label{rem:simeq-not-approx-submeas}
	The implication above does not extend to submeasurements. For example, if
	$A_a^x=0$ for all $a$, then $A_a^x \ot I \simeq_0 I \ot B_a^x$, but
	\[
		\E_{\bx} \sum_a \lVert (A^{\bx}_a \ot I - I \ot B^{\bx}_a)\ket{\psi} \rVert^2
		= \E_{\bx} \sum_a \lVert (I \ot B^{\bx}_a)\ket{\psi} \rVert^2,
	\]
	which is nonzero unless $(I \ot B_a^x)\ket{\psi}=0$ for all $x$ and $a$.
\end{remark}

\section{Consistency from state-dependent distance}
\label{sec:consistency-from-state-dependent-distance}

\begin{theorem}[Closeness of inner products]\label{prop:closeness-of-ip}
	\lean{MIPStarRE.LDT.Preliminaries.closenessOfIP, MIPStarRE.LDT.Preliminaries.closenessOfIPAdjoint}
	\leanok
	\uses{def:approx_delta}
	Let $\{A^x_a\}$, $\{B^x_a\}$, and $\{C^x_{a,b}\}$ be matrices. Suppose that $A^x_a \approx_\gamma B^x_a$ and that for all~$x$,
	\[
		\sum_a \Big(\sum_b C^x_{a,b}\Big) \Big(\sum_b C^x_{a,b}\Big)^\dagger \le I.
	\]
	Then
	\begin{equation}
		\E_{\bx} \sum_{a,b} \bra{\psi} C^{\bx}_{a,b} A^{\bx}_a  \ket{\psi} \approx_{\sqrt{\gamma}} \E_{\bx} \sum_{a,b} \bra{\psi} C^{\bx}_{a,b} B^{\bx}_a  \ket{\psi}. \label{eq:closeness3}
	\end{equation}
	Similarly, suppose that $(A^x_a)^\dagger \approx_\gamma (B^x_a)^\dagger$ and that for all~$x$,
	\[
		\sum_a \Big(\sum_b C^x_{a,b}\Big)^\dagger \Big(\sum_b C^x_{a,b}\Big) \le I.
	\]
	Then
	\begin{equation}
		\E_{\bx} \sum_{a,b} \bra{\psi} A^{\bx}_a  C^{\bx}_{a,b} \ket{\psi} \approx_{\sqrt{\gamma}} \E_{\bx} \sum_{a,b} \bra{\psi} B^{\bx}_a  C^{\bx}_{a,b}  \ket{\psi}. \label{eq:closeness4}
	\end{equation}
\end{theorem}
\begin{proof}
	\leanok
	To prove \eqref{eq:closeness3}, write the difference as
	\[
		\E_{\bx} \sum_a \bra{\psi} \Big(\sum_b C^\bx_{a,b}\Big) (A^\bx_a-B^\bx_a) \ket{\psi}.
	\]
	Cauchy--Schwarz bounds its magnitude by
	\[
		\Big( \E_{\bx} \sum_a \bra{\psi} \Big(\sum_b C^{\bx}_{a,b}\Big)\Big(\sum_b C^{\bx}_{a,b}\Big)^\dagger \ket{\psi} \Big)^{1/2}
		\Big( \E_{\bx} \sum_a \bra{\psi} (A^{\bx}_a - B^{\bx}_a)^\dagger (A^{\bx}_a - B^{\bx}_a) \ket{\psi} \Big)^{1/2},
	\]
	and the two hypotheses bound these factors by $1$ and $\sqrt{\gamma}$ respectively. For \eqref{eq:closeness4}, rewrite the difference as
	\[
		\E_{\bx} \sum_{a,b} \bra{\psi} (C^{\bx}_{a,b})^\dagger \big((A^{\bx}_a)^\dagger-(B^{\bx}_a)^\dagger\big) \ket{\psi},
	\]
	and apply \eqref{eq:closeness3} to the adjoint families.
\end{proof}

\begin{theorem}\label{prop:easy-approx-from-approx-delta}
	\lean{MIPStarRE.LDT.Preliminaries.easyApproxFromApproxDelta}
	\leanok
	\uses{def:submeasurement, def:approx_delta}
	Let $A = \{A^x_a\}$, $B = \{B^x_a\}$, and $C = \{C^x_a\}$ be submeasurements such that $A^x_a \approx_{\delta} B^x_a$. Then
	\begin{equation*}
		\E_{\bx} \sum_a \bra{\psi} A^{\bx}_a  C^{\bx}_a \ket{\psi}
		\approx_{\sqrt{\delta}} \E_{\bx} \sum_a \bra{\psi} B^{\bx}_a  C^{\bx}_a \ket{\psi}.
	\end{equation*}
\end{theorem}
\begin{proof}
	\leanok
	The difference has magnitude
	\[
		\Big| \E_{\bx} \sum_a \bra{\psi} (A^{\bx}_a - B^{\bx}_a) C^{\bx}_a \ket{\psi}\Big|.
	\]
	By Cauchy--Schwarz this is at most
	\[
		\Big(\E_{\bx} \sum_a \bra{\psi} (A^{\bx}_a - B^{\bx}_a)^2 \ket{\psi}\Big)^{1/2}
		\Big(\E_{\bx} \sum_a \bra{\psi} (C^{\bx}_a)^2 \ket{\psi}\Big)^{1/2}.
	\]
	The first factor is at most $\sqrt{\delta}$ by hypothesis, and the second is at most $1$ because $C$ is a submeasurement.
\end{proof}

\begin{theorem}\label{prop:cab-approx-delta}
	\lean{MIPStarRE.LDT.Preliminaries.cabApproxDelta}
	\leanok
	\uses{def:approx_delta}
	Let $\{A^x_a\}$, $\{B^x_a\}$, and $\{C^x_{a,b}\}$ be matrices. Suppose that $A^{x}_a \approx_\delta B^{x}_a$ and that for all $x$ and $a$,
	\[
		\sum_b (C^{x}_{a,b})^\dagger (C^{x}_{a,b}) \leq I.
	\]
	Then
	\[
		C^{x}_{a,b} A^x_{a} \approx_{\delta} C^{x}_{a,b} B^{x}_a.
	\]
\end{theorem}
\begin{proof}
	\leanok
	The error term is
	\[
		\E_{\bx} \sum_{a,b} \bra{\psi}
		(A^{\bx}_a - B^{\bx}_a)^\dagger
		(C^{\bx}_{a,b})^\dagger C^{\bx}_{a,b}
		(A^{\bx}_a - B^{\bx}_a) \ket{\psi}.
	\]
	Using $\sum_b (C^{x}_{a,b})^\dagger (C^{x}_{a,b}) \le I$, this is bounded by
	\[
		\E_{\bx} \sum_a \bra{\psi}
		(A^{\bx}_a - B^{\bx}_a)^\dagger
		(A^{\bx}_a - B^{\bx}_a) \ket{\psi}
		\le \delta.
	\]
\end{proof}

\begin{lemma}[Triangle inequality for vectors squared]\label{prop:triangle-inequality-for-vectors-squared}
	\lean{MIPStarRE.LDT.Preliminaries.triangleInequalityForVectorsSquared}
	\leanok
	For vectors $\ket{\psi_1},\dots,\ket{\psi_k}$,
	\[
		\bigl\lVert \ket{\psi_1}+\cdots+\ket{\psi_k} \bigr\rVert^2 \le k(\lVert\ket{\psi_1}\rVert^2+\cdots+\lVert\ket{\psi_k}\rVert^2).
	\]
\end{lemma}
\begin{proof}
	\leanok
	First,
	\begin{equation}\label{eq:prop-for-real-numbers}
		(x_1+\cdots+x_k)^2 \le k(x_1^2+\cdots+x_k^2)
	\end{equation}
	for all real numbers $x_1,\dots,x_k$. Apply \eqref{eq:prop-for-real-numbers} to $x_i=\lVert\ket{\psi_i}\rVert$ and combine it with the triangle inequality.
\end{proof}

\begin{lemma}[Triangle inequality for $\approx_\delta$]\label{prop:triangle-inequality-for-approx_delta}
	\lean{MIPStarRE.LDT.Preliminaries.triangleInequalityForApproxDelta}
	\leanok
	\uses{def:approx_delta, prop:triangle-inequality-for-vectors-squared}
	Suppose $(A_i)_a^x \approx_{\delta_i} (A_{i+1})_a^x$ for $i=1,\dots,k$. Then
	\[
		(A_1)_a^x \approx_{k(\delta_1+\cdots+\delta_k)} (A_{k+1})_a^x.
	\]
\end{lemma}
\begin{proof}
	Expand $(A_1-A_{k+1})\ket{\psi}$ as a telescoping sum and apply Lemma~\ref{prop:triangle-inequality-for-vectors-squared}.
\end{proof}

\begin{lemma}[Triangle inequality for $\simeq_\delta$]\label{prop:simeq-triangle-inequality}
	\lean{MIPStarRE.LDT.Preliminaries.simeqTriangleInequality}
	\leanok
	\uses{prop:simeq-to-approx, prop:triangle-inequality-for-approx_delta, prop:triangle-sub}
	Suppose $A_a^x \ot I \simeq_\eps I \ot B_a^x$, $C_a^x \ot I \simeq_\delta I \ot B_a^x$, and $C_a^x \ot I \simeq_\gamma I \ot D_a^x$, where all four families are measurements. Then
	\[
		A_a^x \ot I \simeq_{\eps + 2\sqrt{\delta+\gamma}} I \ot D_a^x.
	\]
\end{lemma}
\begin{proof}
	\leanok
	Convert the two consistencies involving $C$ to state-dependent distance, compose them by the triangle inequality, and transfer the result back to consistency by Lemma~\ref{prop:triangle-sub}.
\end{proof}

\begin{lemma}[Data processing for consistency]\label{prop:simeq-data-processing}
	\lean{MIPStarRE.LDT.Preliminaries.simeqDataProcessing}
	\leanok
	\uses{def:post-processing, def:simeq}
	If $A=\{A_a^x\}$ and $B=\{B_a^x\}$ are measurements such that $A_a^x \ot I \simeq_\delta I \ot B_a^x$, and $f$ is a map on the answer set, then
	\[
		A_{[f(a)=b]}^x \ot I \simeq_\delta I \ot B_{[f(a)=b]}^x.
	\]
\end{lemma}
\begin{proof}
	Merging outcome classes can only decrease the off-diagonal mass.
\end{proof}

\begin{lemma}[Consistency controls a submeasurement]\label{prop:cons-sub-meas}
	\lean{MIPStarRE.LDT.Preliminaries.consSubMeas}
	\leanok
	\uses{def:simeq, def:approx_delta, prop:triangle-inequality-for-approx_delta}
	Let $A=\{A_a^x\}$ be a submeasurement and $B=\{B_a^x\}$ a measurement such that $A_a^x \ot I \simeq_\gamma I \ot B_a^x$. Then
	\begin{equation}
		A_a^x \ot I \approx_\gamma A_a^x \ot B_a^x \approx_\gamma A^x \ot B_a^x,
		\label{eq:closeness5}
	\end{equation}
	where $A^x=\sum_a A_a^x$. In particular,
	\[
		A_a^x \ot I \approx_{4\gamma} A^x \ot B_a^x.
	\]
\end{lemma}
\begin{proof}
	For the first approximation in \eqref{eq:closeness5}, bound
	\[
		\E_{\bx} \sum_a \bra{\psi} \bigl(A^{\bx}_a \ot (I-B^{\bx}_a)\bigr)^2 \ket{\psi}
	\]
	by the inconsistency between $A$ and $B$. The second approximation is identical, applied to
	\[
		\E_{\bx} \sum_a \bra{\psi} \bigl((A^\bx-A^\bx_a)\ot B^\bx_a\bigr)^2 \ket{\psi}.
	\]
	Then compose the two bounds with Lemma~\ref{prop:triangle-inequality-for-approx_delta}.
\end{proof}

\begin{lemma}[Switch-sandwich]\label{prop:switch-sandwich}
	\lean{MIPStarRE.LDT.Preliminaries.switchSandwich}
	\leanok
	\uses{def:approx_delta}
	Suppose $A=\{A_a^x\}$ is a projective submeasurement satisfying
	\begin{equation}
		A_a^x \ot I \approx_\delta I \ot A_a^x.
		\label{eq:Aapproxd}
	\end{equation}
	Then for every operator $B$ with $0 \le B \le I$,
	\begin{equation}
		\E_x \sum_a \bra{\psi} A_a^x B A_a^x \ot I \ket{\psi}
		\approx_{2\sqrt\delta}
		\E_x \sum_a \bra{\psi} B \ot A_a^x \ket{\psi}
		\approx_{\sqrt\delta}
		\E_x \sum_a \bra{\psi} BA_a^x \ot I \ket{\psi}.
		\label{eq:switch-sandwich}
	\end{equation}
\end{lemma}
\begin{proof}
	We prove the first approximation in \eqref{eq:switch-sandwich} in two steps. First,
	\begin{equation}
		\E_{\bx} \sum_a \bra{\psi} A^{\bx}_a B A^{\bx}_a \ot I \ket{\psi}
		\approx_{\sqrt{\delta}}
		\E_{\bx} \sum_a \bra{\psi} A^{\bx}_a B \ot A^{\bx}_a \ket{\psi}.
		\label{eq:shift-right-A}
	\end{equation}
	This is obtained by applying Cauchy--Schwarz to the difference and using $B^2 \le B \le I$ together with \eqref{eq:Aapproxd}. For the second step, move the remaining left copy of $A^\bx_a$ across the bipartition in the same way; projectivity then collapses the sandwich and yields the first approximation in \eqref{eq:switch-sandwich}. The second approximation in \eqref{eq:switch-sandwich} is the same argument with the rightmost factor $A_a^\bx$ omitted.
\end{proof}

\section{Strong self-consistency}

\begin{lemma}[Completeness transfer to a projective approximation]\label{prop:completeness-transfer-projective-P}
	\lean{MIPStarRE.LDT.Preliminaries.completenessTransferProjectiveP}
	\leanok
	\uses{def:approx_delta, prop:easy-approx-from-approx-delta}
	Let $A=\{A_a^x\}$ be a submeasurement and $P=\{P_a^x\}$ a projective submeasurement such that $A_a^x \ot I \approx_\eps P_a^x \ot I$. Then
	\[
		\bra{\psi} A \ot I \ket{\psi} \ge \bra{\psi} P \ot I \ket{\psi} - 2\sqrt\eps.
	\]
\end{lemma}
\begin{proof}
	Since $P$ is projective,
	\[
		\bra{\psi} P \ot I \ket{\psi}
		= \E_{\bx} \sum_a \bra{\psi} (P^{\bx}_a)^2 \ot I \ket{\psi}.
	\]
	Apply Proposition~\ref{prop:easy-approx-from-approx-delta} twice to replace $(P^\bx_a)^2$ first by $A^\bx_aP^\bx_a$ and then by $(A^\bx_a)^2$, and conclude with $(A_a^\bx)^2 \le A_a^\bx$.
\end{proof}

\begin{definition}[Strong self-consistency]\label{def:strong-self-consistency}
	\lean{MIPStarRE.LDT.SSCRel, MIPStarRE.LDT.BipartiteSSCRel}
	\leanok
	Let $\ket{\psi}$ be permutation-invariant and let $A=\{A_a^x\}$ be a submeasurement. We say that $A$ is $\delta$-strongly self-consistent when
	\[
		\E_x \sum_a \bra{\psi} A_a^x \ot A_a^x \ket{\psi}
		\ge
		\bra{\psi} A \ot I \ket{\psi} - \delta.
	\]
\end{definition}

\begin{lemma}[Strong self-consistency implies ordinary self-consistency]\label{prop:other-two-notions-of-self-consistency}
	\lean{MIPStarRE.LDT.Preliminaries.otherTwoNotionsOfSelfConsistency}
	\leanok
	\uses{def:strong-self-consistency, def:simeq}
	If $A$ is $\delta$-strongly self-consistent, then $A_a^x \ot I \simeq_\delta I \ot A_a^x$. If $A$ is a measurement, the converse also holds.
\end{lemma}
\begin{proof}
	\leanok
	Rewrite the off-diagonal mass as the difference between the total mass of $A$ and the diagonal overlap in the strong self-consistency inequality.
\end{proof}

\begin{lemma}[Strong self-consistency and state-dependent distance]\label{prop:two-notions-of-self-consistency}
	\lean{MIPStarRE.LDT.Preliminaries.twoNotionsOfSelfConsistency}
	\leanok
	\uses{def:strong-self-consistency, def:approx_delta}
	If $A$ is $\delta$-strongly self-consistent, then
	\[
		A_a^x \ot I \approx_{2\delta} I \ot A_a^x.
	\]
	If $A$ is projective, the converse also holds.
\end{lemma}
\begin{proof}
	Expanding the squared norm gives
	\begin{align}
		      & \E_{\bx} \sum_a \lVert(A^{\bx}_a \otimes I - I \otimes A^{\bx}_a)\ket{\psi}\rVert^2 \nonumber                                                                  \\
		=~    & 2\cdot \Big( \E_{\bx} \sum_a \bra{\psi} (A^{\bx}_a)^2 \otimes I \ket{\psi} -  \E_{\bx} \sum_a \bra{\psi} A^{\bx}_a \otimes A^{\bx}_a \ket{\psi}\Big) \nonumber \\
		\leq~ & 2\cdot \Big( \E_{\bx} \sum_a \bra{\psi} A^{\bx}_a \otimes I \ket{\psi} -  \E_{\bx} \sum_a \bra{\psi} A^{\bx}_a \otimes A^{\bx}_a \ket{\psi}\Big).
		\label{eq:here's-where-projectivity-would-help}
	\end{align}
	The strong self-consistency bound controls the last line by $2\delta$. If $A$ is projective, then \eqref{eq:here's-where-projectivity-would-help} is an equality, which gives the converse.
\end{proof}

\begin{lemma}[Post-processing preserves strong self-consistency]\label{prop:two-notions-of-self-consistency-after-evaluation}
	\lean{MIPStarRE.LDT.Preliminaries.twoNotionsOfSelfConsistencyAfterEvaluation}
	\leanok
	\uses{def:post-processing, def:strong-self-consistency, prop:two-notions-of-self-consistency}
	If $A$ is $\delta$-strongly self-consistent and $f$ is a function on its answer set, then
	\[
		A_{[f(a)=b]}^x \ot I \approx_{2\delta} I \ot A_{[f(a)=b]}^x.
	\]
\end{lemma}
\begin{proof}
	\leanok
	The same computation as in Lemma~\ref{prop:two-notions-of-self-consistency} gives
	\begin{align}
		      & \E_{\bx} \sum_b \lVert(A^{\bx}_{[f(a)=b]} \otimes I - I \otimes A^{\bx}_{[f(a)=b]}) \ket{\psi}\rVert^2 \nonumber                            \\
		\leq~ & 2\cdot \Big( \bra{\psi} A \otimes I \ket{\psi} -  \E_{\bx} \sum_b \bra{\psi} A^{\bx}_{[f(a)=b]} \otimes A^{\bx}_{[f(a)=b]} \ket{\psi}\Big).
		\label{eq:finishing-this-up}
	\end{align}
	The post-processed diagonal term dominates $\E_{\bx} \sum_a \bra{\psi} A^\bx_a \otimes A^\bx_a \ket{\psi}$, so strong self-consistency bounds \eqref{eq:finishing-this-up} by $2\delta$.
\end{proof}

\begin{lemma}[Completeness transfer for a strongly self-consistent approximation]\label{prop:completeness-transfer-self-consistent-A}
	\lean{MIPStarRE.LDT.Preliminaries.completenessTransferSelfConsistentA}
	\leanok
	\uses{def:strong-self-consistency, def:approx_delta, prop:easy-approx-from-approx-delta}
	Let $A$ be a $\delta$-strongly self-consistent submeasurement and let $B$ be a submeasurement such that $A_a^x \ot I \approx_\eps B_a^x \ot I$. Then
	\[
		\bra{\psi} B \ot I \ket{\psi} \ge \bra{\psi} A \ot I \ket{\psi} - \delta - 2\sqrt\eps.
	\]
\end{lemma}
\begin{proof}
	\leanok
	Bound $\bra{\psi} B \ot I \ket{\psi}$ from below by $\E_{\bx}\sum_a \bra{\psi} B^\bx_a \ot B^\bx_a \ket{\psi}$. Then apply Proposition~\ref{prop:easy-approx-from-approx-delta} twice to replace this by $\E_{\bx}\sum_a \bra{\psi} A^\bx_a \ot A^\bx_a \ket{\psi}$, and finish with strong self-consistency of $A$.
\end{proof}

\begin{lemma}[Self-consistency implies data-processing closeness]\label{prop:self-consistency-implies-data-processing}
	\lean{MIPStarRE.LDT.Preliminaries.selfConsistencyImpliesDataProcessing}
	\leanok
	\uses{def:post-processing, def:strong-self-consistency, def:approx_delta, prop:two-notions-of-self-consistency-after-evaluation, prop:completeness-transfer-projective-P, prop:easy-approx-from-approx-delta, prop:triangle-inequality-for-approx_delta}
	Let $A$ be a $\delta$-strongly self-consistent submeasurement and let $P$ be a projective submeasurement such that $P_a^x \ot I \approx_\eps A_a^x \ot I$. Then for every function $f$,
	\[
		P_{[f(a)=b]}^x \ot I \approx_{8\delta + 8\sqrt\eps} A_{[f(a)=b]}^x \ot I.
	\]
\end{lemma}
\begin{proof}
	\leanok
	First prove
	\begin{equation}
		P^x_{[f(a)=b]} \ot I \approx_{2\delta + 4\sqrt{\eps}} I \ot A_{[f(a)=b]}^x.
		\label{eq:what-we-want-to-prove-but-on-wrong-side}
	\end{equation}
	Expanding the squared norm and using projectivity gives
	\begin{align}
		      & \E_{\bx} \sum_b \lVert(P^{\bx}_{[f(a)=b]} \ot I - I \ot A_{[f(a)=b]}^{\bx}) \ket{\psi} \rVert^2 \nonumber \\
		\leq~ & \bra{\psi} P \ot I \ket{\psi} + \bra{\psi} I \ot A \ket{\psi}
		- 2\cdot \E_{\bx} \sum_b \bra{\psi}  P^{\bx}_{[f(a)=b]} \ot A^{\bx}_{[f(a)=b]} \ket{\psi}.
		\label{eq:gonna-handle-third-term}
	\end{align}
	Proposition~\ref{prop:completeness-transfer-projective-P} bounds the first term by $\bra{\psi} A \ot I \ket{\psi}+2\sqrt{\eps}$, and Proposition~\ref{prop:easy-approx-from-approx-delta} together with strong self-consistency bounds the third term below by $\bra{\psi} A \ot I \ket{\psi}-\delta-\sqrt{\eps}$. Substituting into \eqref{eq:gonna-handle-third-term} yields \eqref{eq:what-we-want-to-prove-but-on-wrong-side}. Now combine \eqref{eq:what-we-want-to-prove-but-on-wrong-side} with
	\[
		A^x_{[f(a)=b]} \ot I \approx_{2\delta} I \ot A^x_{[f(a)=b]}
	\]
	from Lemma~\ref{prop:two-notions-of-self-consistency-after-evaluation}, and conclude by Lemma~\ref{prop:triangle-inequality-for-approx_delta}.
\end{proof}

\begin{lemma}[Self-consistency moves one copy to the same side]\label{lem:self-consistency-same-side-square}
	\lean{MIPStarRE.LDT.Preliminaries.bipartiteSSCSquaredMass}
	\leanok
	\uses{def:strong-self-consistency}
	Let $A=\{A_a\}$ be a $\zeta$-strongly self-consistent submeasurement. Then
	\[
		\sum_a \bra{\psi} A_a^2 \ot I \ket{\psi}
		\ge
		\sum_a \bra{\psi} A_a \ot I \ket{\psi} - \zeta.
	\]
\end{lemma}
\begin{proof}
	Apply Cauchy--Schwarz to the two-sided overlap $\sum_a \bra{\psi} A_a \ot A_a \ket{\psi}$ and compare it with the strong self-consistency lower bound.
\end{proof}

\begin{lemma}[Bounding the missing mass of a close submeasurement]\label{lem:completion-missing-mass-bound}
	\lean{MIPStarRE.LDT.Preliminaries.completionMissingMassBound}
	\uses{def:strong-self-consistency, def:approx_delta, lem:self-consistency-same-side-square}
	Let $A=\{A_a\}$ be a measurement that is $\zeta$-strongly self-consistent, and let $B=\{B_a\}$ be a submeasurement such that $A_a \ot I \approx_\delta B_a \ot I$. Writing $B=\sum_a B_a$, one has
	\[
		\lVert (I-B) \ot I\ket{\psi} \rVert^2 \le 2\sqrt\delta + \zeta.
	\]
\end{lemma}
\begin{proof}
	Since $0 \le I-B \le I$, the left-hand side is at most $1-\bra{\psi} B \ot I \ket{\psi}$. Compare $\sum_a \bra{\psi} B_a^2 \ot I \ket{\psi}$ with $\sum_a \bra{\psi} A_a B_a \ot I \ket{\psi}$ and then with $\sum_a \bra{\psi} A_a^2 \ot I \ket{\psi}$ using the $\approx_\delta$ hypothesis, and finish with Lemma~\ref{lem:self-consistency-same-side-square}.
\end{proof}

\begin{lemma}[Completing a close submeasurement to a full measurement]\label{prop:completing-to-measurement}
	\lean{MIPStarRE.LDT.Preliminaries.completingToMeasurement}
	\leanok
	\uses{def:measurement-completion, def:strong-self-consistency, def:approx_delta, prop:triangle-inequality-for-vectors-squared, prop:easy-approx-from-approx-delta, prop:cool-prop}
	Let $A=\{A_a\}$ be a measurement that is $\zeta$-strongly self-consistent, let $B=\{B_a\}$ be a submeasurement such that $A_a \ot I \approx_\delta B_a \ot I$, and let $C$ be the measurement obtained by adding the missing mass $I-B$ to one distinguished answer $a^*$. Then
	\[
		A_a \ot I \approx_{2\delta + 4\sqrt\delta + 2\zeta} C_a \ot I.
	\]
\end{lemma}

% prop:cool-prop from the paper is the same result as lem:self-consistency-same-side-square above.
% We keep the label for cross-reference fidelity.
\begin{theorem}\label{prop:cool-prop}
	\lean{MIPStarRE.LDT.Preliminaries.bipartiteSSCSquaredMass}
	\leanok
	\uses{lem:self-consistency-same-side-square}
	This is the same statement as Lemma~\ref{lem:self-consistency-same-side-square}: for a $\zeta$-strongly self-consistent submeasurement $A$,
	\[
		\sum_a \bra{\psi} (A_a)^2 \ot I \ket{\psi} \geq \sum_a \bra{\psi} A_a \ot I \ket{\psi} -\zeta.
	\]
\end{theorem}
\begin{proof}
	\leanok
	\uses{lem:self-consistency-same-side-square}
	Immediate from Lemma~\ref{lem:self-consistency-same-side-square}.
\end{proof}

\begin{proof}[Proof of \ref{prop:completing-to-measurement}]
	By Lemma~\ref{prop:triangle-inequality-for-vectors-squared},
	\[
		\sum_a \lVert(A_a-C_a)\ot I \ket{\psi} \rVert^2
		\leq 2\delta + 2 \cdot \lVert(I-B)\ot I \ket{\psi}\rVert^2.
	\]
	It remains to bound the residual term. Since $0 \le I-B \le I$,
	\begin{align}
		\lVert(I-B) \ot I \ket{\psi}\rVert^2
		 & \leq 1 -\sum_a \bra{\psi} B_a^2 \ot I \ket{\psi}.
		\label{eq:to-return-later-whatevs}
	\end{align}
	By Proposition~\ref{prop:easy-approx-from-approx-delta},
	\[
		\sum_a \bra{\psi} B_a^2 \ot I \ket{\psi}
		\approx_{\sqrt{\delta}} \sum_a \bra{\psi} (A_a B_a)\ot I \ket{\psi}
		\approx_{\sqrt{\delta}} \sum_a \bra{\psi} A_a^2\ot I \ket{\psi}.
	\]
	Proposition~\ref{prop:cool-prop} gives
	\[
		\sum_a \bra{\psi} A_a^2\ot I \ket{\psi} \geq \bra{\psi} A \ot I \ket{\psi} - \zeta = 1-\zeta.
	\]
	Hence \eqref{eq:to-return-later-whatevs} implies
	\[
		\lVert(I-B) \ot I \ket{\psi}\rVert^2 \le 2\sqrt\delta + \zeta.
	\]
	Substituting this into the first bound proves the claim.
\end{proof}
